% =============================================================================
%  1-Introduccion.tex — Introducción del registro de avances
%  Describe el repositorio, su propósito y el periodo cubierto por el reporte.
% =============================================================================

\section{Introducción}
\label{sec:introduccion}

% ── INSTRUCCIONES ─────────────────────────────────────────────────────────────
% Describe brevemente:
%   1. ¿Qué es el repositorio \NombreRepo y cuál es su propósito?
%   2. ¿Qué periodo cubre este reporte (de \VerRef\ a \VerActual)?
%   3. ¿En qué contexto se enmarca (materia, semestre, investigación)?
%
% Las variables \NombreRepo, \VerRef y \VerActual se definen en main.tex.
% Puedes usarlas directamente en el texto:
%   "El repositorio \texttt{\NombreRepo} contiene..."
% ─────────────────────────────────────────────────────────────────────────────

El presente documento registra los avances del repositorio \texttt{\NombreRepo}
durante el periodo comprendido entre la versión \texttt{\VerRef} y la versión
\texttt{\VerActual}. El seguimiento se realiza a través del historial de
\termino{commits} y \termino{tags} del sistema de control de versiones
(\termino{Git}), lo que permite rastrear con precisión cada cambio estructural
o de contenido realizado sobre el documento.

% ── Añade aquí tu descripción del proyecto y contexto ─────────────────────────
% TODO: redactar introducción del registro.
